\centerline{\bf \Huge Геометрические неравенства}

\begin{flushright}
        \scriptsize{Я очень устал\\
                    Мне хочется спать\\
                    Мне хочется знать...}
\end{flushright}

\problem{1}
Внутри треугольника $A B C$ выбрана произвольная точка $M$. Докажите, что:

(a) $B M+C M<A B+A C$;

(b) сумма расстояний от $M$ до вершин треугольника больше полупериметра, но меньше периметра треугольника;

(c) сумма расстояний от $M$ до вершин треугольника не превосходит суммы двух наибольших сторон.

\problem{2}
На основании $A C$ равнобедренного треугольника $A B C$ выбрали точку $D$, а на его продолжении за вершину $C$ - точку $E$, причем $A D=C E$ Докажите, что $B D+B E> A B+B C$.

\problem{3}
Точки $D$ и $E$ делят сторону $A C$ треугольника $A B C$ на три равные части. Докажите, что $B D+B E<A B+B C$.

\problem{4}
На стороне $A D$ выпуклого четырёхугольника $A B C D$ нашлась такая точка $M$, что $C M$ и $B M$ параллельны $A B$ и $C D$ соответственно. Докажите, что $S_{A B C D} \geqslant 3 S_{B C M}$.

\problem{5}
На стороне $B E$ правильного треугольника $A B E$ вне его построен ромб $B C D E$. Отрезки $A C$ и $B D$ пересекаются в точке $F$. Докажите, что $A F<B D$.

\problem{6}
Пусть $O$ - центр описанной окружности треугольника $A B C$. На сторонах $A B$ и $B C$ выбраны точки $M$ и $N$ соответственно, причём $2 \angle M O N=\angle A O C$. Докажите, что периметр треугольника $M B N$ не меньше стороны $A C$.

\problem{7}
Внутри выпуклого четырёхугольника $A B C D$, в котором $A B=C D$, выбрана точка $P$ таким образом, что сумма углов $P B A$ и $P C D$ равна $180^{\circ}$. Докажите, что $P B+P C<A D$.

\problem{8}
Центр описанной окружности треугольника лежит на вписанной окружности. Докажите, что отношение наибольшей стороны треугольника к наименьшей меньше 2.